%%%%%%%%%%%%%%%%%%%%%%%%%%%%%%%%%%%%%%%%%%%%%%%%%%%%%%%%%%%%%%%%%%%%%%%%%%%%%%%%%
%Beamer-Präsentationsvorlage
%%%%%%%%%%%%%%%%%%%%%%%%%%%%%%%%%%%%%%%%%%%%%%%%%%%%%%%%%%%%%%%%%%%%%%%%%%%%%%%%%

\documentclass[xcolor=dvipsnames]{beamer}


%\usepackage{xcolor}
\usetheme{default}
% deutsche Silbentrennung
\usepackage[ngerman]{babel}
% Deutsche Umlaute
\usepackage[utf8]{inputenc}
% Einbinden von Bildern
\usepackage{graphicx}
% Unterstützung von Hyperlinks
\usepackage{hyperref}

% Mathematische Pakete
\usepackage{amsmath}
\usepackage{amssymb}
\usepackage{amsfonts}
\usepackage{amsthm}
\usepackage{latexsym}
\usepackage{colortbl}
\usepackage{listings}

% Abweichung von Vorlage
% um Graphen zu plotten
\usepackage{tikz}
\usetikzlibrary{shapes, arrows, automata, chains, positioning, shadows, matrix}



\usepackage{tikzscale}
\usepackage{tikz-uml}\usepackage{tabularx}


\makeatletter
\newcommand{\trickbeamer}{%
	\advance\beamer@slideinframe by-1%
}%
\makeatother



\def \variableXYZ{Ueberschrift1}
\def \variableXYZ{Ueberschrift1}
\newcommand{\neueFolie}[1]{\frametitle{#1}\section{#1}}


\definecolor{rgray}{rgb}{0.95,0.95,0.95}
\definecolor{darkgreen}{rgb}{0.0,0.5,0.0}
\lstset{
	language=[LaTeX]TeX, 
	captionpos=b, 
	frame=lines, 
	basicstyle=\ttfamily, 
	keywordstyle=\color{blue}, 
	commentstyle=\color{darkgreen}, 
	stringstyle=\color{red}, 
	%numbers=left, 
	%numberstyle=\tiny, 
	%numbersep=5pt,
	breaklines=true, 
	showstringspaces=false, 
	emph={double,bool,int,unsigned,char,true,false,void},
	emphstyle=\color{blue},
	emph={Assert,Test},
	emphstyle=\color{red},
	emph={[2]\using,\#define,\#ifdef,\#endif}, 
	emphstyle={[2]\color{blue}},
	backgroundcolor=\color{rgray}
}

\setbeamercolor{frametitle}{fg=black}
\setbeamertemplate{navigation symbols}{}%remove navigation symbols
\setbeamercolor{itemize item}{bg=gray, fg=gray} 
\setbeamercolor{itemize subitem}{fg=gray}
\setbeamercolor{section in toc}{fg=darkgray}
\setbeamercolor{subsection in toc shaded}{fg=gray}
\setbeamercolor{title}{fg=black}
\usepackage{animate}
%\setbeamertemplate{frametitle}
%{
%\includegraphics[width=3cm]{Gestaltung/oth_logo.png}\hfill
%%\parbox{0.6\textwidth}{\color{black}
%%\vspace*{-1cm}{
%\insertframetitle
%%}
%%}
%\hspace*{-0.3cm}\color{black}\rule{20cm}{0.2mm}
%}

\setbeamertemplate{frametitle}
{% 
%\vspace{-0.165ex} 
\begin{beamercolorbox}[wd=\paperwidth,dp=1ex, ht=4ex, sep=0.5ex, colsep*=0pt]{frametitle}% 
    \usebeamerfont{frametitle}   \strut 
    \raisebox{-0.55ex}[0pt][-\ht\strutbox ]{
    	\includegraphics[width=1.7cm]{Gestaltung/Logo_Hochschule_Ansbach.png}
    	}  
    \hfill  \insertframetitle \hspace{0.4cm}
    \end{beamercolorbox}
    \vspace*{-0.5cm}
    \hspace*{-2cm}
\color{black}\rule{15cm}{0.2mm}               
}%


%\setbeamertemplate{headline}{
%\begin{beamercolorbox}{frametitle}


%\begin{flushright}

%\end{flushright}
%}} 
%\end{beamercolorbox}%

%}
\setbeamertemplate{footline}
{
  \leavevmode%
  \vspace*{0.2cm}
     \hspace*{-2cm}
	\color{black}\rule{15cm}{0.2mm}   
  \hbox{%
  \begin{beamercolorbox}[wd=1.15\paperwidth,ht=1.5ex,dp=1ex, center]{author in head/foot}%
    \usebeamerfont{author in head/foot}\color{black}\insertshortauthor \hfill
\insertshorttitle  \hfill    Seite \insertframenumber /\inserttotalframenumber
\vspace*{0.02cm}
  \end{beamercolorbox}}%
}

\setbeamerfont{tiny structure}{size=\tiny}


\defbeamertemplate*{title page}{customized}[1][]
{
\begin{minipage}{0.8\linewidth}
  \usebeamerfont{title}\inserttitle\par
  \usebeamerfont{subtitle}\usebeamercolor[fg]{subtitle}\insertsubtitle\par
  \bigskip
  \usebeamerfont{author}\insertauthor\par
  \usebeamerfont{institute}\insertinstitute\par
  \usebeamerfont{date}\insertdate\par
  %\usebeamercolor[fg]{titlegraphic}\inserttitlegraphic
\end{minipage} \hfill
\begin{minipage}{0.15\linewidth}
  \vspace*{0pt}
  \usebeamercolor[fg]{titlegraphic}\inserttitlegraphic
\end{minipage}
  
  
}

%%%%%%%%%%%%%%%%%%%%%%%%%%%%%%%%%%%%%%%%%%%%%%%%%%%%%%%%%%%%%%%%%%%%%%%%%%%%%%%%%
%Hier werden die Daten für die erste Folie eingetragen
%%%%%%%%%%%%%%%%%%%%%%%%%%%%%%%%%%%%%%%%%%%%%%%%%%%%%%%%%%%%%%%%%%%%%%%%%%%%%%%%%
\title{ UML Tools}
\subtitle{Open-Source Software for the Working Scientist \\}
\institute{Master of Applied Research\\ Hochschule Ansbach\\}
\author{Marco Obermeier}
\date{\today}
%\date{12. Januar 2022}

\begin{document} 
	\begin{frame}
		\frametitle{~}
		\maketitle 
	\end{frame}
	%%%%%%%%%%%%%%%%%%%%%%%%%%%%%%%%%%%%%%%%%%%%%%%%%%%%%%%%%%%%%%%%%%%%%%%%%%%%%%%%%
	%Einleitung
	%%%%%%%%%%%%%%%%%%%%%%%%%%%%%%%%%%%%%%%%%%%%%%%%%%%%%%%%%%%%%%%%%%%%%%%%%%%%%%%%%
	%\begin{frame}
	%	\frametitle{Inhaltsverzeichnis}
	%	\tableofcontents
	%\end{frame}
	%%%%%%%%%%%%%%%%%%%%%%%%%%%%%%%%%%%%%%%%%%%%%%%%%%%%%%%%%%%%%%%%%%%%%%%%%%%%%%%%%
	%Einleitung
	%%%%%%%%%%%%%%%%%%%%%%%%%%%%%%%%%%%%%%%%%%%%%%%%%%%%%%%%%%%%%%%%%%%%%%%%%%%%%%%%%
\begin{frame}
	\frametitle{Übersicht und Entwurf}
	\begin{columns}
		\begin{column}{0.70\textwidth}
			\begin{minipage}{1.0\textwidth}	
				\resizebox{\columnwidth}{!}{%
					\begin{tikzpicture}
						
					
						\umlbasiccomponent[x=0, y=0, fill=yellow!20]{MySQL} 
						\umlbasiccomponent[x=4, y=0,fill=yellow!20]{Sauter Vision Center}
						\umlbasiccomponent[x=8.5, y=0,fill=yellow!20]{Wetterstation}
						\umlbasiccomponent[x=12, y=0,fill=yellow!20]{Sensorik}
						\umlbasiccomponent[x=6.5, y=-4, fill=cyan!20]{Datenkollektor} 
	
						%\umldatabase[x=5, y=-6, class=InfluxDB]{Zeitreihendatenbank}
			
						\umlHVassemblyconnector[interface=historische Daten, anchor1=180,  first arm=1.5cm]{Datenkollektor}{MySQL}
						\umlVHVassemblyconnector[interface={Gebäudedaten}, anchor1=125, anchor2=-135]{Datenkollektor}{Sauter Vision Center}
						\umlVHVassemblyconnector[interface={Wetterdaten}, anchor1=65, anchor2=-45]{Datenkollektor}{Wetterstation}
						\umlHVassemblyconnector[interface=Leistungsdaten, anchor1=0,  first arm=1.5cm]{Datenkollektor}{Sensorik}
						%\umlVHVassemblyconnector[interface=3D-Beschleunigung, anchor1=0,  first arm=3.5cm]{MySQL}{Datenkollektor}
						%\umlassemblyconnector[interface=Winkelgeschwidigkeit, anchor1=0]{Flugsteuerung}{Gyrosensor}
						%\umlHVHassemblyconnector[interface=PWM-Signal, anchor1=-62,  first arm=3.5cm]{Flugsteuerung}{Motoren}
		
						%\umlVHVassemblyconnector[interface=Lagewinkel]{Regelung}{Sensorauswertung} 
						%\umlVHVassemblyconnector[interface=Motordrehzahl]{Motorsteuerung}{Regelung} 
						
						
					\end{tikzpicture}
				}
			\end{minipage}
		\end{column}
		\begin{column}{0.30\textwidth}
			\begin{block}{Kategeroien der Komponenten}
				\begin{itemize}
					\item blau$\equiv$ Hardware
					\item gelb $\equiv$ Software
					\item rot	$\equiv$ Echtzeit-kritisch
				\end{itemize}
				
				
			\end{block}
		\end{column}
	\end{columns}
\end{frame}
	
	
	
	
	
	
\end{document}